\documentclass[11pt]{amsart}

% Encoding and fonts
\usepackage[T1]{fontenc}
\usepackage[utf8]{inputenc}
\usepackage{lmodern}

% Math packages
\usepackage{amsmath}
\usepackage{amssymb}
\usepackage{amsfonts}
\usepackage{amsthm}
\usepackage{mathtools}

% Layout and utilities
\usepackage{geometry}
\usepackage{enumitem}
\usepackage{booktabs}
\usepackage{longtable}
\usepackage[hidelinks]{hyperref}

\geometry{margin=1in}

\numberwithin{equation}{section}
\setcounter{tocdepth}{3}
\setcounter{secnumdepth}{3}

% Theorem-like environments
\theoremstyle{plain}
\newtheorem{theorem}{Theorem}[section]
\newtheorem{proposition}[theorem]{Proposition}
\newtheorem{lemma}[theorem]{Lemma}
\newtheorem{corollary}[theorem]{Corollary}
\newtheorem{conjecture}[theorem]{Conjecture}

\theoremstyle{definition}
\newtheorem{definition}[theorem]{Definition}
\newtheorem{experiment}[theorem]{Experiment}
\newtheorem{heuristic}[theorem]{Heuristic}
\newtheorem{question}[theorem]{Question}

\theoremstyle{remark}
\newtheorem{remark}[theorem]{Remark}
\newtheorem{warning}[theorem]{Warning}

\newcommand{\dd}{\,\mathrm d}
\newcommand{\e}{\mathrm e}
\newcommand{\ii}{\mathrm i}
\newcommand{\dist}{\operatorname{dist}}
\newcommand{\Li}{\operatorname{Li}}
\newcommand{\ZP}{\mathfrak Z}
\newcommand{\N}{\mathbb N}
\newcommand{\R}{\mathbb R}
\newcommand{\C}{\mathbb C}

\title[Marked $z$-Primes and a Spectral Lattice]{
Marked $z$-Primes Attached to the Nontrivial Zeros of the Riemann Zeta Function:\
Exact Quantization, Prime-Power Resonances, Asymptotic Bounds, and a Theoretical Localization Lattice
}

\author{}
\date{}

\begin{document}

\begin{abstract}
Let $p_k$ be the $k$-th ordinary prime and let $\gamma_k>0$ be the
$k$-th positive ordinate of a nontrivial zero of the Riemann zeta
function, counted with multiplicity.  We define the marked $z$-prime
attached to $\gamma_k$ by
\begin{equation}
 n_k=\left\lceil\frac{\gamma_k\log p_k}{\pi}\right\rceil,
 \qquad
 z_k=\exp\left(\frac{\pi n_k}{\gamma_k}\right),
 \qquad
 \mathfrak z_k=(n_k,z_k).
 \label{eq:abstract_definition}
\end{equation}
The integer mark is essential: the pair $(n_k,z_k)$ recovers the ordinate
exactly through $\gamma_k=\pi n_k/\log z_k$, whereas $z_k$ alone retains
only a phase class.  The construction is local, canonical, and need not be
monotone in the second coordinate.  We prove exact quantization identities,
a bijection between the positive zero ordinates and the marked system,
sharp one-cell bounds, the asymptotics $n_k\sim2k$, $z_k\sim p_k$, and
\begin{equation}
 z_k-p_k
 =\left(\frac{\varepsilon_k}{2}+o(1)\right)(\log k)^2,
 \qquad 0\leq\varepsilon_k<1.
 \label{eq:abstract_displacement}
\end{equation}
Consequently the upper envelope satisfies
$\limsup (z_k-p_k)/(\log k)^2\leq1/2$, and the marked counting function has
the same first-order density as the ordinary primes.

The logarithmic coordinate $\log z_k$ is the first node above $\log p_k$
in the variable lattice $(\pi/\gamma_k)\mathbb Z$.  This produces exact
formulae for the detuning of the Mangoldt frequencies $r\log p_k$ from
signed resonant nodes.  In particular,
\begin{equation}
 \e^{\ii\gamma_k\log p_k^r}
 =(-1)^{rn_k}\e^{-\ii\pi r\varepsilon_k},
 \label{eq:abstract_prime_power_phase}
\end{equation}
and the distance to the nearest signed resonance is
$\pi\|r\varepsilon_k\|_{\R/\mathbb Z}$.  Dirichlet approximation therefore
produces arbitrarily accurate resonant prime powers, while the factor
$p_k^{-r/2}\log p_k$ from the explicit formula suppresses high powers.
Finally, we formulate a bivariate marked Dirichlet transform and an inverse
lattice program in which the pairs $(n,z)$ serve as exact coordinates for
zero ordinates.  All inverse-selection statements beyond the exact transform
identities are isolated as open analytic problems.
\end{abstract}

\maketitle
\tableofcontents

\section{Introduction}
\label{sec:introduction}

The nontrivial zeros of the Riemann zeta function are written
\begin{equation}
 \rho_k=\beta_k+\ii\gamma_k,
 \qquad
 0<\gamma_1\leq\gamma_2\leq\cdots,
 \label{eq:ordered_zeros}
\end{equation}
where ordinates are counted with multiplicity.  The classical
Riemann--von Mangoldt formula controls their counting function, while the
prime number theorem controls the ordered primes $p_k$; standard references
are \cite{Titchmarsh1986,Ingham1932}.  The present construction combines the
$k$-th ordinate with the $k$-th prime through an exact logarithmic phase
quantization.

The central object is not an unmarked real number.  It is the pair
\begin{equation}
 \mathfrak z_k=(n_k,z_k),
 \label{eq:introduction_marked_pair}
\end{equation}
where $n_k$ records the lifted phase and $z_k$ records the associated
positive scale.  This distinction is structural.  The identity
\begin{equation}
 \gamma_k\log z_k=\pi n_k
 \label{eq:introduction_quantization}
\end{equation}
can be inverted only after $n_k$ is retained.

The construction has three purposes.  First, it produces a canonical marked
encoding of every positive zero ordinate.  Second, it places the logarithm of
an ordinary prime inside a single explicitly controlled phase cell.  Third,
it defines a theoretical lattice in which a candidate pair $(n,z)$ determines
a candidate ordinate $\pi n/\log z$ exactly.

\begin{remark}
\label{rem:status_convention}
All statements concerning the definition, exact inversion, asymptotic bounds,
counting density, and prime-power detuning are proved rigorously below.
Statements concerning an independent analytic rule that would select the
correct marked points without first inserting the ordinates are explicitly
labelled as questions, heuristics, or conjectures.
\end{remark}

\section{Classical input}
\label{sec:classical_input}

\subsection{Zero and prime clocks}
\label{subsec:zero_prime_clocks}

Let
\begin{equation}
 N(T)=\#\{k:0<\gamma_k\leq T\}
 \label{eq:zero_counting_function}
\end{equation}
with the usual midpoint convention at an ordinate when needed.  The
Riemann--von Mangoldt formula gives
\begin{equation}
 N(T)
 =\frac{T}{2\pi}\log\left(\frac{T}{2\pi}\right)
 -\frac{T}{2\pi}
 +O(\log T),
 \qquad T\to\infty,
 \label{eq:riemann_von_mangoldt}
\end{equation}
see \cite[Chapter~9]{Titchmarsh1986}.  Inverting the leading term yields
\begin{equation}
 \gamma_k\sim\frac{2\pi k}{\log k}.
 \label{eq:gamma_asymptotic}
\end{equation}

The prime number theorem and its inverse form give
\begin{align}
 \pi(x)&\sim\frac{x}{\log x},
 \label{eq:pnt_pi}
 \\
 p_k&\sim k\log k.
 \label{eq:nth_prime_asymptotic}
\end{align}
The inverse asymptotic and its refinements are classical; see
\cite{Ingham1932,AriasToulisse2013}.

\begin{lemma}
\label{lem:prime_zero_ratio}
As $k\to\infty$,
\begin{align}
 \frac{p_k}{\gamma_k}
 &\sim\frac{(\log k)^2}{2\pi},
 \label{eq:prime_zero_ratio}
 \\
 \frac{\pi p_k}{\gamma_k}
 &\sim\frac12(\log k)^2.
 \label{eq:scaled_prime_zero_ratio}
\end{align}
\end{lemma}

\begin{proof}
Divide \eqref{eq:nth_prime_asymptotic} by
\eqref{eq:gamma_asymptotic}.  Multiplication by $\pi$ gives
\eqref{eq:scaled_prime_zero_ratio}.
\end{proof}

\subsection{Mangoldt frequencies}
\label{subsec:mangoldt_frequencies}

The von Mangoldt function is supported on prime powers:
\begin{equation}
 \Lambda(q)=
 \begin{cases}
 \log p,&q=p^r,\quad r\geq1,\\
 0,&\text{otherwise}.
 \end{cases}
 \label{eq:mangoldt_definition}
\end{equation}
For $\Re s>1$,
\begin{equation}
 -\frac{\zeta'}{\zeta}(s)
 =\sum_{q=2}^{\infty}\frac{\Lambda(q)}{q^s}.
 \label{eq:log_derivative_dirichlet_series}
\end{equation}
In the Guinand--Weil explicit formula, the arithmetic frequencies occur at
$\pm\log q$ with $q=p^r$ and coefficient scale
$\Lambda(q)/\sqrt q$ for the critical-line normalization; see
\cite{Guinand1948,Weil1952}.

\section{Canonical marked $z$-primes}
\label{sec:canonical_definition}

\begin{definition}[Local marked $z$-prime]
\label{def:local_marked_zprime}
For each $k\geq1$, define
\begin{align}
 a_k
 &:=\frac{\gamma_k\log p_k}{\pi},
 \label{eq:ak_definition}
 \\
 n_k
 &:=\lceil a_k\rceil,
 \label{eq:nk_definition}
 \\
 \varepsilon_k
 &:=n_k-a_k,
 \qquad 0\leq\varepsilon_k<1,
 \label{eq:epsilon_definition}
 \\
 z_k
 &:=\exp\left(\frac{\pi n_k}{\gamma_k}\right),
 \label{eq:zk_definition}
 \\
 \mathfrak z_k
 &:=(n_k,z_k).
 \label{eq:marked_zprime_definition}
\end{align}
The marked system is
\begin{equation}
 \ZP:=\{\mathfrak z_k:k\geq1\}.
 \label{eq:marked_system_definition}
\end{equation}
\end{definition}

\begin{remark}
\label{rem:ceiling_choice}
The ceiling convention chooses the first resonant lattice node at or above
$p_k$.  A nearest-node convention would reduce the maximal logarithmic error
by a factor of two but would permit $z_k<p_k$.  The one-sided convention is
retained because it gives a canonical upper cell and an exact one-sided
comparison with $p_k$.
\end{remark}

\begin{proposition}[Exact quantization]
\label{prop:exact_quantization}
For every $k\geq1$,
\begin{align}
 \gamma_k\log z_k
 &=\pi n_k,
 \label{eq:exact_quantization}
 \\
 \e^{\ii\gamma_k\log z_k}
 &=(-1)^{n_k},
 \label{eq:exact_signed_phase}
 \\
 z_k
 &=p_k\exp\left(\frac{\pi\varepsilon_k}{\gamma_k}\right).
 \label{eq:exact_zprime_factorization}
\end{align}
\end{proposition}

\begin{proof}
Equation \eqref{eq:exact_quantization} follows directly from
\eqref{eq:zk_definition}.  Exponentiation gives
\eqref{eq:exact_signed_phase}.  Finally,
\eqref{eq:epsilon_definition} and \eqref{eq:ak_definition} give
\begin{equation}
 \frac{\pi n_k}{\gamma_k}
 =\log p_k+\frac{\pi\varepsilon_k}{\gamma_k},
 \label{eq:proof_log_factorization}
\end{equation}
which exponentiates to \eqref{eq:exact_zprime_factorization}.
\end{proof}

\begin{theorem}[Exact recovery and marked bijection]
\label{thm:exact_recovery_bijection}
The map
\begin{equation}
 \Phi:\{\gamma_k:k\geq1\}\longrightarrow\ZP,
 \qquad
 \Phi(\gamma_k)=(n_k,z_k),
 \label{eq:Phi_definition}
\end{equation}
is a bijection when the zero ordinates are counted with their indexing
multiplicity.  Its inverse is
\begin{equation}
 \Phi^{-1}(n,z)=\frac{\pi n}{\log z}.
 \label{eq:Phi_inverse}
\end{equation}
\end{theorem}

\begin{proof}
Surjectivity holds by the definition of $\ZP$.  If
$(n_j,z_j)=(n_k,z_k)$, then \eqref{eq:exact_quantization} gives
\begin{equation}
 \gamma_j
 =\frac{\pi n_j}{\log z_j}
 =\frac{\pi n_k}{\log z_k}
 =\gamma_k.
 \label{eq:bijection_injectivity_identity}
\end{equation}
Thus the marked data determine the ordinate exactly.  With multiplicities
retained at the index level, each indexed zero corresponds to one indexed
marked pair.
\end{proof}

\begin{warning}
\label{warn:unmarked_projection}
The projection $(n_k,z_k)\mapsto z_k$ can lose information.  Distinct
ordinates could in principle produce the same positive coordinate with
different integer marks.  The exact object is therefore the marked pair, not
the unmarked real number.
\end{warning}

\section{The variable logarithmic lattice}
\label{sec:variable_lattice}

\begin{definition}[Lattice attached to an ordinate]
\label{def:lattice_attached_ordinate}
For $\gamma>0$, define the signed logarithmic resonance lattice
\begin{equation}
 \mathcal L(\gamma)
 :=\left\{\frac{\pi n}{\gamma}:n\in\mathbb Z\right\}
 \subset\mathbb R,
 \label{eq:log_lattice_definition}
\end{equation}
and its positive exponential image
\begin{equation}
 \mathcal E(\gamma)
 :=\left\{\exp\left(\frac{\pi n}{\gamma}\right):n\in\mathbb Z\right\}.
 \label{eq:exponential_lattice_definition}
\end{equation}
\end{definition}

The mesh in logarithmic coordinate is
\begin{equation}
 h(\gamma)=\frac{\pi}{\gamma}.
 \label{eq:lattice_mesh}
\end{equation}
For the $k$-th zero, $\log z_k$ is the first node of
$\mathcal L(\gamma_k)$ at or above $\log p_k$.

\begin{proposition}[One-cell characterization]
\label{prop:one_cell_characterization}
The marked $z$-prime is uniquely characterized by
\begin{align}
 \log z_k&\in\mathcal L(\gamma_k),
 \label{eq:one_cell_membership}
 \\
 \log p_k&\leq\log z_k<\log p_k+\frac{\pi}{\gamma_k}.
 \label{eq:one_cell_interval}
\end{align}
Equivalently,
\begin{equation}
 \frac{\pi(n_k-1)}{\gamma_k}
 <\log p_k
 \leq\frac{\pi n_k}{\gamma_k}.
 \label{eq:cell_index_inequality}
\end{equation}
\end{proposition}

\begin{proof}
The ceiling relation
$n_k-1<a_k\leq n_k$ is equivalent to
\eqref{eq:cell_index_inequality}.  The right endpoint is $\log z_k$ by
\eqref{eq:zk_definition}, and subtracting $\log p_k$ gives
\eqref{eq:one_cell_interval}.  There is exactly one first lattice node in
that half-open cell.
\end{proof}

\begin{remark}
\label{rem:not_single_fixed_lattice}
The construction is a family of lattices rather than a single fixed lattice:
the mesh $\pi/\gamma_k$ varies with $k$.  The marked sequence selects one
node from each member of this family.  In the $(n,\log z)$ plane, every
selected point determines the ordinate through the level curve
$\gamma=\pi n/\log z$.
\end{remark}

\section{Exact and asymptotic bounds}
\label{sec:bounds}

\subsection{Sharp one-cell inequalities}
\label{subsec:sharp_one_cell_inequalities}

\begin{theorem}[Exact local bounds]
\label{thm:exact_local_bounds}
For every $k\geq1$,
\begin{align}
 0
 &\leq\log\frac{z_k}{p_k}
 <\frac{\pi}{\gamma_k},
 \label{eq:exact_log_bound}
 \\
 p_k
 &\leq z_k
 <p_k\e^{\pi/\gamma_k},
 \label{eq:exact_multiplicative_bound}
 \\
 0
 &\leq z_k-p_k
 <p_k\left(\e^{\pi/\gamma_k}-1\right).
 \label{eq:exact_additive_bound}
\end{align}
Moreover,
\begin{align}
 \frac{\pi(n_k-1)}{\log p_k}
 &<\gamma_k
 \leq\frac{\pi n_k}{\log p_k},
 \label{eq:gamma_bounds_from_n_p}
 \\
 \gamma_k
 &=\frac{\pi n_k}{\log z_k}.
 \label{eq:gamma_exact_from_mark}
\end{align}
\end{theorem}

\begin{proof}
Equation \eqref{eq:exact_log_bound} is
\eqref{eq:one_cell_interval}.  Exponentiation yields
\eqref{eq:exact_multiplicative_bound}; subtracting $p_k$ gives
\eqref{eq:exact_additive_bound}.  Rearranging
\eqref{eq:cell_index_inequality} yields
\eqref{eq:gamma_bounds_from_n_p}.  Equation
\eqref{eq:gamma_exact_from_mark} is the exact inversion formula.
\end{proof}

\begin{corollary}[Bounds using only $n_k$ and $p_k$]
\label{cor:n_p_bounds_z}
For $n_k>1$,
\begin{equation}
 p_k\leq z_k<p_k^{n_k/(n_k-1)}.
 \label{eq:z_bound_n_p}
\end{equation}
\end{corollary}

\begin{proof}
The upper bound in \eqref{eq:gamma_bounds_from_n_p} gives
$\log z_k\geq\log p_k$.  The lower bound gives
\begin{equation}
 \log z_k
 =\frac{\pi n_k}{\gamma_k}
 <\frac{n_k}{n_k-1}\log p_k,
 \label{eq:proof_n_p_upper}
\end{equation}
which proves \eqref{eq:z_bound_n_p}.
\end{proof}

\subsection{Asymptotic characterization}
\label{subsec:asymptotic_characterization}

\begin{theorem}[Asymptotic mark, scale, and displacement]
\label{thm:asymptotic_characterization}
As $k\to\infty$,
\begin{align}
 \frac{n_k}{k}&\longrightarrow2,
 \label{eq:nk_over_k_limit}
 \\
 \frac{z_k}{p_k}&\longrightarrow1,
 \label{eq:zk_over_pk_limit}
 \\
 \frac{z_k}{k\log k}&\longrightarrow1,
 \label{eq:zk_scale_limit}
 \\
 z_k-p_k
 &=\left(\frac{\varepsilon_k}{2}+o(1)\right)(\log k)^2.
 \label{eq:precise_displacement_asymptotic}
\end{align}
The $o(1)$ in \eqref{eq:precise_displacement_asymptotic} is uniform for
$0\leq\varepsilon_k<1$.
\end{theorem}

\begin{proof}
From \eqref{eq:nk_definition} and \eqref{eq:epsilon_definition},
\begin{equation}
 n_k
 =\frac{\gamma_k\log p_k}{\pi}+\varepsilon_k.
 \label{eq:nk_exact_epsilon}
\end{equation}
Using \eqref{eq:gamma_asymptotic},
$\log p_k\sim\log k$, and $0\leq\varepsilon_k<1$ gives
\eqref{eq:nk_over_k_limit}.

By \eqref{eq:exact_zprime_factorization},
\begin{equation}
 \log\frac{z_k}{p_k}
 =\frac{\pi\varepsilon_k}{\gamma_k}
 =O\left(\frac{\log k}{k}\right),
 \label{eq:log_ratio_rate}
\end{equation}
which proves \eqref{eq:zk_over_pk_limit}.  Combining with
$p_k\sim k\log k$ gives \eqref{eq:zk_scale_limit}.

Finally, uniformly for $0\leq\varepsilon_k<1$,
\begin{equation}
 \e^{\pi\varepsilon_k/\gamma_k}-1
 =\frac{\pi\varepsilon_k}{\gamma_k}
 +O\left(\frac{1}{\gamma_k^2}\right).
 \label{eq:uniform_exponential_expansion}
\end{equation}
Multiply by $p_k$ and apply Lemma~\ref{lem:prime_zero_ratio}.  Since
$p_k/\gamma_k^2=o((\log k)^2)$, equation
\eqref{eq:precise_displacement_asymptotic} follows.
\end{proof}

\begin{corollary}[Maximal asymptotic envelope]
\label{cor:maximal_asymptotic_envelope}
The displacement satisfies
\begin{align}
 z_k-p_k&=O((\log k)^2),
 \label{eq:displacement_big_o}
 \\
 0&\leq\limsup_{k\to\infty}
 \frac{z_k-p_k}{(\log k)^2}
 \leq\frac12.
 \label{eq:displacement_limsup_bound}
\end{align}
More precisely,
\begin{equation}
 \frac{z_k-p_k}{(\log k)^2}
 -\frac{\varepsilon_k}{2}
 \longrightarrow0.
 \label{eq:displacement_epsilon_limit}
\end{equation}
\end{corollary}

\begin{proof}
Apply \eqref{eq:precise_displacement_asymptotic} and
$0\leq\varepsilon_k<1$.
\end{proof}

\begin{remark}
\label{rem:upper_envelope_not_attained}
The constant $1/2$ is the universal asymptotic upper envelope obtained from
the one-cell construction.  Equality of the limsup with $1/2$ would require
values of $\varepsilon_k$ approaching $1$ along a subsequence.  No such
distribution statement is assumed here.
\end{remark}

\begin{corollary}[Scaled logarithmic displacement]
\label{cor:scaled_log_displacement}
One has
\begin{equation}
 \frac{k}{\log k}\log\frac{z_k}{p_k}
 =\frac{\varepsilon_k}{2}+o(1).
 \label{eq:scaled_log_displacement_limit}
\end{equation}
Consequently,
\begin{equation}
 0\leq\limsup_{k\to\infty}
 \frac{k}{\log k}\log\frac{z_k}{p_k}
 \leq\frac12.
 \label{eq:scaled_log_limsup}
\end{equation}
\end{corollary}

\begin{proof}
Use \eqref{eq:exact_log_bound}, the exact identity
$\log(z_k/p_k)=\pi\varepsilon_k/\gamma_k$, and
\eqref{eq:gamma_asymptotic}.
\end{proof}

\section{Counting density and generalized-prime interpretation}
\label{sec:counting_density}

\begin{definition}[Marked and projected counting functions]
\label{def:counting_functions}
Define
\begin{align}
 \Pi_{\ZP}(x)
 &:=\#\{k:z_k\leq x\},
 \label{eq:projected_counting_function}
 \\
 \Pi_{\ZP}(x;A)
 &:=\#\{k:z_k\leq x,\ n_k\in A\},
 \qquad A\subseteq\N.
 \label{eq:marked_counting_function}
\end{align}
The first count includes multiplicity in the projected coordinate.
\end{definition}

\begin{theorem}[Prime-density preservation]
\label{thm:prime_density_preservation}
The projected counting function satisfies
\begin{equation}
 \Pi_{\ZP}(x)\sim\pi(x)\sim\frac{x}{\log x}.
 \label{eq:zprime_counting_asymptotic}
\end{equation}
\end{theorem}

\begin{proof}
Since $z_k\geq p_k$, one has
\begin{equation}
 \Pi_{\ZP}(x)\leq\pi(x).
 \label{eq:counting_upper}
\end{equation}
Fix $\delta>0$.  By \eqref{eq:zk_over_pk_limit}, for all sufficiently
large $k$,
\begin{equation}
 z_k\leq(1+\delta)p_k.
 \label{eq:eventual_multiplicative_delta}
\end{equation}
Therefore, apart from finitely many indices,
$p_k\leq x/(1+\delta)$ implies $z_k\leq x$, and hence
\begin{equation}
 \pi\left(\frac{x}{1+\delta}\right)-O_\delta(1)
 \leq\Pi_{\ZP}(x)
 \leq\pi(x).
 \label{eq:counting_squeeze}
\end{equation}
Divide by $x/\log x$, use the prime number theorem, and let
$\delta\downarrow0$.
\end{proof}

\begin{remark}
\label{rem:ordering_not_required}
The proof does not require the sequence $(z_k)$ to be increasing.  Counting
is performed on the indexed multiset.  If one wishes to place the projected
coordinates into the standard formalism of Beurling generalized primes, they
may be rearranged nondecreasingly and counted with multiplicity; the Euler
product in its half-plane of absolute convergence is independent of that
rearrangement.  The general framework of generalized primes originates in
\cite{Beurling1937}; see also \cite{Diamond1969}.
\end{remark}

\begin{definition}[Projected Beurling-type zeta function]
\label{def:projected_zeta_function}
For $\Re s>1$, define
\begin{equation}
 \zeta_{\ZP}^{\mathrm{proj}}(s)
 :=\prod_{k=1}^{\infty}(1-z_k^{-s})^{-1}.
 \label{eq:projected_euler_product}
\end{equation}
\end{definition}

\begin{proposition}[Absolute convergence]
\label{prop:projected_absolute_convergence}
The product \eqref{eq:projected_euler_product} converges absolutely and
locally uniformly for $\Re s>1$.
\end{proposition}

\begin{proof}
Since $z_k\geq p_k$,
\begin{equation}
 \sum_{k=1}^{\infty}z_k^{-\sigma}
 \leq\sum_{k=1}^{\infty}p_k^{-\sigma}<\infty,
 \qquad \sigma>1.
 \label{eq:absolute_convergence_comparison}
\end{equation}
The standard Euler-product convergence criterion applies.
\end{proof}

\begin{warning}
\label{warn:projected_product_forgets_marks}
The projected product \eqref{eq:projected_euler_product} forgets $n_k$ and
therefore does not retain the exact inverse map \eqref{eq:Phi_inverse}.  The
marked transform introduced in Section~\ref{sec:marked_transforms} is the
appropriate object for zero recovery.
\end{warning}

\section{Mangoldt-frequency resonances}
\label{sec:mangoldt_resonances}

\subsection{The matched prime}
\label{subsec:matched_prime}

The ordinary prime $p_k$ lies below the exact resonant node $z_k$ in the same
logarithmic cell.  Its phase is therefore determined by the cell remainder.

\begin{theorem}[Exact matched-prime phase]
\label{thm:exact_matched_prime_phase}
For every $k\geq1$,
\begin{align}
 \gamma_k\log p_k
 &=\pi(n_k-\varepsilon_k),
 \label{eq:matched_prime_phase_lift}
 \\
 \e^{\ii\gamma_k\log p_k}
 &=(-1)^{n_k}\e^{-\ii\pi\varepsilon_k}.
 \label{eq:matched_prime_phase}
\end{align}
The angular distance of $\gamma_k\log p_k$ from the signed lattice
$\pi\mathbb Z$ is
\begin{equation}
 \dist(\gamma_k\log p_k,\pi\mathbb Z)
 =\pi\min\{\varepsilon_k,1-\varepsilon_k\}.
 \label{eq:matched_prime_signed_distance}
\end{equation}
\end{theorem}

\begin{proof}
Equation \eqref{eq:matched_prime_phase_lift} is the definition of
$\varepsilon_k$.  Exponentiation yields
\eqref{eq:matched_prime_phase}.  The distance formula follows because
$n_k$ is an integer and $0\leq\varepsilon_k<1$.
\end{proof}

\begin{remark}
\label{rem:metric_closeness_phase_detuning}
The metric closeness $z_k/p_k\to1$ does not force
$\e^{\ii\gamma_k\log p_k}$ to approach $\pm1$.  Multiplication by
$\gamma_k$ expands the logarithmic displacement back to the full cell
coordinate $\pi\varepsilon_k$.  Thus $z_k$ is close to $p_k$ in scale, while
the residual phase within the cell can range across an interval of length
$\pi$.
\end{remark}

\subsection{Prime powers}
\label{subsec:prime_powers}

\begin{theorem}[Exact prime-power detuning]
\label{thm:exact_prime_power_detuning}
For every $r\geq1$,
\begin{align}
 \gamma_k\log p_k^r
 &=\pi r(n_k-\varepsilon_k),
 \label{eq:prime_power_lift}
 \\
 \e^{\ii\gamma_k\log p_k^r}
 &=(-1)^{rn_k}\e^{-\ii\pi r\varepsilon_k},
 \label{eq:prime_power_phase}
 \\
 \dist(\gamma_k\log p_k^r,\pi\mathbb Z)
 &=\pi\|r\varepsilon_k\|_{\R/\mathbb Z},
 \label{eq:prime_power_signed_detuning}
\end{align}
where
\begin{equation}
 \|x\|_{\R/\mathbb Z}:=\min_{m\in\mathbb Z}|x-m|.
 \label{eq:circle_distance_definition}
\end{equation}
\end{theorem}

\begin{proof}
Multiply \eqref{eq:matched_prime_phase_lift} by $r$.  The phase identity
follows by exponentiation, and the distance identity follows after removing
the integer $rn_k$.
\end{proof}

\begin{corollary}[Finite resonant-power approximation]
\label{cor:finite_resonant_power_approximation}
For every real $\varepsilon_k$ and every integer $R\geq1$, there exists an
integer $1\leq r\leq R$ such that
\begin{equation}
 \dist(\gamma_k\log p_k^r,\pi\mathbb Z)
 \leq\frac{\pi}{R+1}.
 \label{eq:dirichlet_resonant_power_bound}
\end{equation}
\end{corollary}

\begin{proof}
Apply the pigeonhole form of Dirichlet approximation to the $R+1$ points
$0,\varepsilon_k,2\varepsilon_k,\ldots,R\varepsilon_k$ modulo one.  Two lie
within distance $1/(R+1)$, so for some $1\leq r\leq R$,
$\|r\varepsilon_k\|_{\R/\mathbb Z}\leq1/(R+1)$.  Apply
\eqref{eq:prime_power_signed_detuning}.
\end{proof}

\begin{proposition}[Exact resonance criterion]
\label{prop:exact_resonance_criterion}
For a fixed $k$ and $r\geq1$, the prime power $p_k^r$ is exactly signed
resonant with $\gamma_k$ if and only if
\begin{equation}
 r\varepsilon_k\in\mathbb Z.
 \label{eq:exact_signed_resonance_criterion}
\end{equation}
It is exactly positive resonant if and only if
\begin{equation}
 r(n_k-\varepsilon_k)\in2\mathbb Z.
 \label{eq:exact_positive_resonance_criterion}
\end{equation}
\end{proposition}

\begin{proof}
The first statement follows from
\eqref{eq:prime_power_signed_detuning}.  The second follows from
\eqref{eq:prime_power_lift} and the condition that the lifted phase be an
even multiple of $\pi$.
\end{proof}

\subsection{Interaction with explicit-formula weights}
\label{subsec:explicit_formula_weights}

At the prime power $q=p_k^r$, the critical-line arithmetic weight is
\begin{equation}
 W_{k,r}
 :=\frac{\Lambda(p_k^r)}{\sqrt{p_k^r}}
 =\frac{\log p_k}{p_k^{r/2}}.
 \label{eq:mangoldt_prime_power_weight}
\end{equation}

\begin{proposition}[Resonance--amplitude tradeoff]
\label{prop:resonance_amplitude_tradeoff}
For every $R\geq1$, there exists $1\leq r\leq R$ satisfying
\eqref{eq:dirichlet_resonant_power_bound}, but its weight obeys
\begin{equation}
 W_{k,r}\leq\frac{\log p_k}{\sqrt{p_k}},
 \label{eq:weight_upper_all_r}
\end{equation}
and for $r\geq2$,
\begin{equation}
 W_{k,r}\leq\frac{\log p_k}{p_k}.
 \label{eq:weight_upper_proper_power}
\end{equation}
Thus arbitrarily fine phase approximation obtained by increasing $r$ is
accompanied by exponential decay in $r$ at fixed $p_k$.
\end{proposition}

\begin{proof}
The existence of $r$ is Corollary~\ref{cor:finite_resonant_power_approximation}.
The bounds follow directly from \eqref{eq:mangoldt_prime_power_weight}.
\end{proof}

\begin{definition}[Weighted resonance score]
\label{def:weighted_resonance_score}
For $R\geq1$, define
\begin{equation}
 \mathcal R_k(R)
 :=\max_{1\leq r\leq R}
 \frac{\log p_k}{p_k^{r/2}}
 \left|\cos\bigl(\pi r\varepsilon_k\bigr)\right|.
 \label{eq:weighted_resonance_score}
\end{equation}
\end{definition}

\begin{remark}
\label{rem:score_interpretation}
The score \eqref{eq:weighted_resonance_score} separates two mechanisms:
phase recurrence of $r\varepsilon_k$ modulo one and Mangoldt-weight decay.
The first is Diophantine; the second is arithmetic and strongly favors the
smallest exponents.  This makes the matched prime and the first few powers the
natural range for any analytic test based on the explicit formula.
\end{remark}

\subsection{General prime-power frequencies near a marked node}
\label{subsec:general_frequencies_near_node}

\begin{proposition}[Detuning from an arbitrary prime power]
\label{prop:arbitrary_prime_power_detuning}
Let $q=p^r$ be any prime power.  Define
\begin{equation}
 \Delta_k(q):=\gamma_k\log\left(\frac{q}{z_k}\right).
 \label{eq:general_detuning_definition}
\end{equation}
Then
\begin{equation}
 \e^{\ii\gamma_k\log q}
 =(-1)^{n_k}\e^{\ii\Delta_k(q)}.
 \label{eq:general_detuning_phase}
\end{equation}
Moreover, if
\begin{equation}
 \left|\log\frac{q}{z_k}\right|\leq\frac{\eta}{\gamma_k},
 \label{eq:general_resonance_window}
\end{equation}
then
\begin{equation}
 \left|\e^{\ii\gamma_k\log q}-(-1)^{n_k}\right|
 \leq\eta.
 \label{eq:general_phase_distance_bound}
\end{equation}
\end{proposition}

\begin{proof}
Use $\gamma_k\log z_k=\pi n_k$ to obtain
\eqref{eq:general_detuning_phase}.  Under
\eqref{eq:general_resonance_window}, one has
$|\Delta_k(q)|\leq\eta$, and
$|\e^{\ii x}-1|\leq|x|$ gives
\eqref{eq:general_phase_distance_bound}.
\end{proof}

\section{Marked analytic transforms}
\label{sec:marked_transforms}

\begin{definition}[Marked atomic measure]
\label{def:marked_atomic_measure}
Define the measure on $\N\times(1,\infty)$
\begin{equation}
 \dd\Pi_{\ZP}(n,x)
 :=\sum_{k\geq1}\delta_{n_k}(n)\,\delta_{z_k}(x).
 \label{eq:marked_atomic_measure}
\end{equation}
\end{definition}

\begin{definition}[Marked Dirichlet transform]
\label{def:marked_dirichlet_transform}
For $\Re s>1$ and $|u|\leq1$, define
\begin{align}
 D_{\ZP}(s,u)
 &:=\sum_{k\geq1}u^{n_k}z_k^{-s},
 \label{eq:marked_dirichlet_transform}
 \\
 &=\int_{\N\times(1,\infty)}u^n x^{-s}\,\dd\Pi_{\ZP}(n,x).
 \label{eq:marked_dirichlet_transform_integral}
\end{align}
\end{definition}

\begin{proposition}[Absolute convergence of the marked transform]
\label{prop:marked_transform_convergence}
The series \eqref{eq:marked_dirichlet_transform} converges absolutely and
locally uniformly for $\Re s>1$ and $|u|\leq1$.
\end{proposition}

\begin{proof}
Since $|u|^{n_k}\leq1$ and $z_k\geq p_k$,
\begin{equation}
 \sum_{k\geq1}|u|^{n_k}z_k^{-\sigma}
 \leq\sum_{k\geq1}p_k^{-\sigma}<\infty,
 \qquad \sigma>1.
 \label{eq:marked_transform_majorant}
\end{equation}
\end{proof}

\begin{definition}[Marked Euler transform]
\label{def:marked_euler_transform}
For $\Re s>1$ and $|u|<1$, define
\begin{equation}
 \zeta_{\ZP}(s,u)
 :=\prod_{k\geq1}\left(1-u^{n_k}z_k^{-s}\right)^{-1}.
 \label{eq:marked_euler_transform}
\end{equation}
\end{definition}

\begin{proposition}[Logarithmic expansion]
\label{prop:marked_logarithmic_expansion}
For $\Re s>1$ and $|u|<1$,
\begin{equation}
 \log\zeta_{\ZP}(s,u)
 =\sum_{m\geq1}\frac1m
 \sum_{k\geq1}u^{mn_k}z_k^{-ms}.
 \label{eq:marked_logarithmic_expansion}
\end{equation}
The double series converges absolutely and locally uniformly.
\end{proposition}

\begin{proof}
Expand $-\log(1-w)=\sum_{m\geq1}w^m/m$ with
$w=u^{n_k}z_k^{-s}$.  Absolute convergence follows from
$|u|<1$, $z_k\geq2$, and Proposition~\ref{prop:marked_transform_convergence}.
\end{proof}

\begin{proposition}[Coefficient decomposition]
\label{prop:coefficient_decomposition}
For $|u|<1$,
\begin{equation}
 D_{\ZP}(s,u)=\sum_{n\geq1}A_n(s)u^n,
 \label{eq:marked_power_series}
\end{equation}
where
\begin{equation}
 A_n(s)=\sum_{k:n_k=n}z_k^{-s}.
 \label{eq:marked_coefficient_dirichlet_series}
\end{equation}
Thus the coefficient indexed by $n$ retains precisely the projected scales
carrying that phase lift.
\end{proposition}

\begin{proof}
Group the absolutely convergent series
\eqref{eq:marked_dirichlet_transform} by equal values of $n_k$.
\end{proof}

\section{Theoretical localization of zero ordinates}
\label{sec:theoretical_localization}

\begin{definition}[Inverse lattice map]
\label{def:inverse_lattice_map}
For $n\in\N$ and $z>1$, define
\begin{equation}
 \Gamma(n,z):=\frac{\pi n}{\log z}.
 \label{eq:inverse_lattice_map}
\end{equation}
\end{definition}

\begin{proposition}[Exact localization on the marked support]
\label{prop:exact_localization_support}
For every atom $(n_k,z_k)$ of \eqref{eq:marked_atomic_measure},
\begin{equation}
 \Gamma(n_k,z_k)=\gamma_k.
 \label{eq:exact_localization_identity}
\end{equation}
\end{proposition}

\begin{proof}
This is \eqref{eq:gamma_exact_from_mark}.
\end{proof}

\begin{proposition}[Stability of the inverse coordinate]
\label{prop:inverse_coordinate_stability}
Let $n\geq1$, $z>1$, and perturb the scale to $z\e^\delta$.  Then
\begin{equation}
 \Gamma(n,z\e^\delta)-\Gamma(n,z)
 =-\frac{\pi n\delta}{\log z\,(\log z+\delta)}.
 \label{eq:exact_inverse_perturbation}
\end{equation}
If $|\delta|\leq\tfrac12\log z$, then
\begin{equation}
 \left|\Gamma(n,z\e^\delta)-\Gamma(n,z)\right|
 \leq\frac{2\pi n}{(\log z)^2}|\delta|.
 \label{eq:inverse_stability_bound}
\end{equation}
\end{proposition}

\begin{proof}
Subtract the two rational expressions
$\pi n/(\log z+\delta)$ and $\pi n/\log z$.  The stated denominator bound
gives \eqref{eq:inverse_stability_bound}.
\end{proof}

\begin{corollary}[Stability near the $k$-th marked point]
\label{cor:stability_near_kth_point}
At $(n_k,z_k)$, a logarithmic error $\delta$ produces, to first order,
\begin{equation}
 \Gamma(n_k,z_k\e^\delta)-\gamma_k
 =-\frac{\gamma_k}{\log z_k}\delta+O\left(
 \frac{\gamma_k}{(\log z_k)^2}\delta^2\right).
 \label{eq:local_inverse_linearization}
\end{equation}
Since $\gamma_k/\log z_k\sim2\pi k/(\log k)^2$, localization of
$\gamma_k$ to bounded absolute error requires logarithmic scale accuracy of
order at most $(\log k)^2/k$.
\end{corollary}

\begin{proof}
Expand \eqref{eq:exact_inverse_perturbation} at $\delta=0$ and use
\eqref{eq:gamma_asymptotic} and \eqref{eq:zk_scale_limit}.
\end{proof}

\begin{heuristic}[Marked-transform localization]
\label{heur:marked_transform_localization}
Suppose an analytic continuation or a regularized inversion procedure recovers
the atomic support of $\dd\Pi_{\ZP}$ from $D_{\ZP}(s,u)$ or
$\zeta_{\ZP}(s,u)$.  Coefficient extraction in $u$ separates the integer
marks, Mellin inversion in $s$ recovers the corresponding positive scales,
and \eqref{eq:inverse_lattice_map} then returns the ordinates.  This is a
heuristic program until an independent construction of the marked transform
is available that does not insert the ordinates into
\eqref{eq:nk_definition}--\eqref{eq:zk_definition}.
\end{heuristic}

\begin{question}
\label{question:intrinsic_selector}
Is there a functional equation, extremal principle, trace identity, or
regularized Euler relation whose solutions select the same marked support
$\{(n_k,z_k)\}$ from the ordinary prime data and the integer lattice alone?
\end{question}

\begin{question}
\label{question:fixed_point_operator}
Can one define an operator on approximate ordinate sequences by the cycle
\begin{equation}
 \{t_k\}
 \longmapsto
 \{(n_k(t),z_k(t))\}
 \longmapsto
 \text{zeros of a regularized marked transform}
 \label{eq:fixed_point_cycle}
\end{equation}
for which the actual ordinates form a locally attracting fixed point?
\end{question}

\section{Comparison with ordinary prime frequencies}
\label{sec:comparison_prime_frequencies}

\begin{proposition}[Frequency-cell width]
\label{prop:frequency_cell_width}
The difference between the marked logarithmic frequency and the matched prime
frequency is
\begin{equation}
 0\leq\log z_k-\log p_k<\frac{\pi}{\gamma_k}
 \sim\frac{\log k}{2k}.
 \label{eq:frequency_cell_width_asymptotic}
\end{equation}
Thus each $\log p_k$ lies within one vanishing cell below an exact signed
resonance for $\gamma_k$.
\end{proposition}

\begin{proof}
The exact inequality is \eqref{eq:exact_log_bound}.  The asymptotic cell
width follows from \eqref{eq:gamma_asymptotic}.
\end{proof}

\begin{proposition}[First-order replacement of prime Dirichlet atoms]
\label{prop:first_order_dirichlet_replacement}
Fix a compact set $K\subset\{s\in\C:\Re s>0\}$.  Uniformly for $s\in K$,
\begin{equation}
 z_k^{-s}-p_k^{-s}
 =-s\frac{\pi\varepsilon_k}{\gamma_k}p_k^{-s}
 +O_K\left(\frac{p_k^{-\Re s}}{\gamma_k^2}\right).
 \label{eq:dirichlet_atom_expansion}
\end{equation}
\end{proposition}

\begin{proof}
Use \eqref{eq:exact_zprime_factorization} to write
\begin{equation}
 z_k^{-s}
 =p_k^{-s}\exp\left(-s\frac{\pi\varepsilon_k}{\gamma_k}\right),
 \label{eq:dirichlet_atom_factorization}
\end{equation}
and apply the Taylor expansion of the exponential uniformly on $K$.
\end{proof}

\begin{corollary}[Relative prime-series convergence]
\label{cor:relative_prime_series_convergence}
For every $\sigma>0$,
\begin{equation}
 \sum_{k\geq1}|z_k^{-s}-p_k^{-s}|
 \label{eq:relative_prime_series}
\end{equation}
converges locally uniformly in $\Re s\geq\sigma$.
\end{corollary}

\begin{proof}
By \eqref{eq:dirichlet_atom_expansion}, it is enough to sum
$p_k^{-\sigma}/\gamma_k$.  The classical asymptotics give
\begin{equation}
 \frac{p_k^{-\sigma}}{\gamma_k}
 \ll_\sigma
 \frac{(\log k)^{1-\sigma}}{k^{1+\sigma}},
 \label{eq:relative_series_majorant}
\end{equation}
which is summable for every $\sigma>0$.
\end{proof}

\begin{theorem}[Holomorphic relative Euler factor]
\label{thm:holomorphic_relative_euler_factor}
For $\Re s>1$, set
\begin{equation}
 H(s):=\frac{\zeta_{\ZP}^{\mathrm{proj}}(s)}{\zeta(s)}.
 \label{eq:relative_euler_factor_definition}
\end{equation}
Then $H$ extends to a holomorphic, nonvanishing function on
$\Re s>0$ through
\begin{equation}
 \log H(s)
 =\sum_{m\geq1}\frac1m
 \sum_{k\geq1}\left(z_k^{-ms}-p_k^{-ms}\right).
 \label{eq:relative_logarithm_series}
\end{equation}
Consequently, wherever $\zeta$ is meromorphic in $\Re s>0$,
\begin{equation}
 \zeta_{\ZP}^{\mathrm{proj}}(s)=H(s)\zeta(s)
 \label{eq:relative_factorization}
\end{equation}
provides the corresponding meromorphic continuation of the projected
$z$-prime Euler product.
\end{theorem}

\begin{proof}
For $m=1$, local uniform convergence on every half-plane
$\Re s\geq\sigma>0$ follows from
Corollary~\ref{cor:relative_prime_series_convergence}.  For $m\geq2$,
\begin{equation}
 |z_k^{-ms}-p_k^{-ms}|
 \leq z_k^{-m\sigma}+p_k^{-m\sigma}
 \leq2p_k^{-m\sigma}.
 \label{eq:higher_m_majorant}
\end{equation}
For fixed $\sigma>0$, finitely many $m$ with $m\sigma\leq1$ are handled by
the same difference estimate, and the remaining tail converges absolutely by
comparison with a geometric series in $p_k^{-\sigma}$.  Thus
\eqref{eq:relative_logarithm_series} is holomorphic on $\Re s>0$.
Exponentiation gives a holomorphic nonvanishing $H$.  On $\Re s>1$, the
Euler-product logarithms show that it equals the quotient
\eqref{eq:relative_euler_factor_definition}; analytic continuation gives
\eqref{eq:relative_factorization}.
\end{proof}

\begin{remark}
\label{rem:relative_factorization_meaning}
The factorization \eqref{eq:relative_factorization} is a rigorous analytic
consequence of the small displacement.  It shows that the projected Euler
system differs from the ordinary prime system by a regular nonvanishing
factor in $\Re s>0$.  The marked transform contains additional information
not present in this projection.
\end{remark}

\section{Index geometry and finite candidate windows}
\label{sec:index_geometry}

\begin{proposition}[Asymptotic diagonal]
\label{prop:asymptotic_diagonal}
The marked points lie asymptotically on the diagonal $n=2k$:
\begin{equation}
 n_k=2k+o(k).
 \label{eq:asymptotic_diagonal}
\end{equation}
Moreover,
\begin{equation}
 \log z_k=\log k+\log\log k+o(1).
 \label{eq:log_z_asymptotic}
\end{equation}
\end{proposition}

\begin{proof}
Equation \eqref{eq:asymptotic_diagonal} is equivalent to
\eqref{eq:nk_over_k_limit}.  Equation \eqref{eq:log_z_asymptotic} follows
from $z_k\sim k\log k$.
\end{proof}

\begin{heuristic}[Finite search strip]
\label{heur:finite_search_strip}
For a large index $k$, the asymptotic relation $n_k=2k+o(k)$ restricts the
integer mark to a narrow relative strip around $2k$, while
$z_k=p_k+O((\log k)^2)$ restricts the scale to a polylogarithmic interval
above $p_k$.  A theoretical search may therefore be organized over candidate
pairs
\begin{equation}
 n\in[2k-E_k,2k+E_k]\cap\mathbb Z,
 \qquad
 z\in[p_k,p_k+C(\log k)^2],
 \label{eq:finite_candidate_strip}
\end{equation}
with $E_k=o(k)$ and any fixed $C>1/2$ for sufficiently large $k$.  An
independent selector is still required to choose the actual pair inside this
strip.
\end{heuristic}

\section{Alternative conventions}
\label{sec:alternative_conventions}

\subsection{Nearest-node convention}
\label{subsec:nearest_node}

\begin{definition}[Nearest marked $z$-prime]
\label{def:nearest_marked_zprime}
Define
\begin{equation}
 n_k^{\mathrm{near}}
 :=\left\lfloor\frac{\gamma_k\log p_k}{\pi}+\frac12\right\rfloor,
 \qquad
 z_k^{\mathrm{near}}
 :=\exp\left(\frac{\pi n_k^{\mathrm{near}}}{\gamma_k}\right).
 \label{eq:nearest_definition}
\end{equation}
\end{definition}

\begin{proposition}[Nearest-node error]
\label{prop:nearest_node_error}
The nearest-node convention satisfies
\begin{equation}
 \left|\log\frac{z_k^{\mathrm{near}}}{p_k}\right|
 \leq\frac{\pi}{2\gamma_k}
 \label{eq:nearest_log_bound}
\end{equation}
and therefore
\begin{equation}
 |z_k^{\mathrm{near}}-p_k|
 \leq\left(\frac14+o(1)\right)(\log k)^2.
 \label{eq:nearest_additive_envelope}
\end{equation}
\end{proposition}

\begin{proof}
Nearest-integer rounding gives an error at most $1/2$ in the lifted cell
coordinate.  The additive asymptotic follows as in
Theorem~\ref{thm:asymptotic_characterization}.
\end{proof}

\begin{remark}
\label{rem:choice_comparison}
The nearest convention minimizes metric displacement, while the ceiling
convention preserves $z_k\geq p_k$ and gives a canonical upper node.  Both
retain exact recovery when the integer mark is stored.  The one-sided
convention is preferable when comparing counting functions by monotone
inequalities; the nearest convention is preferable when minimizing local
perturbation.
\end{remark}

\subsection{Positive-phase convention}
\label{subsec:positive_phase}

\begin{definition}[Positive-phase lattice]
\label{def:positive_phase_lattice}
Define
\begin{equation}
 n_k^{+}
 :=\left\lceil\frac{\gamma_k\log p_k}{2\pi}\right\rceil,
 \qquad
 z_k^{+}
 :=\exp\left(\frac{2\pi n_k^{+}}{\gamma_k}\right).
 \label{eq:positive_phase_definition}
\end{equation}
\end{definition}

\begin{proposition}[Positive-phase cost]
\label{prop:positive_phase_cost}
One has
\begin{align}
 \e^{\ii\gamma_k\log z_k^{+}}&=1,
 \label{eq:positive_phase_identity}
 \\
 0\leq\log\frac{z_k^{+}}{p_k}&<\frac{2\pi}{\gamma_k},
 \label{eq:positive_phase_log_bound}
 \\
 z_k^{+}-p_k&\leq(1+o(1))(\log k)^2.
 \label{eq:positive_phase_additive_envelope}
\end{align}
\end{proposition}

\begin{proof}
The proof repeats the ceiling argument with lattice mesh $2\pi/\gamma_k$.
Lemma~\ref{lem:prime_zero_ratio} gives the final envelope.
\end{proof}

\begin{remark}
\label{rem:signed_preferred}
The signed lattice has half the mesh and therefore half the maximal one-sided
displacement.  Retaining the parity of $n_k$ records whether the exact node
has phase $+1$ or $-1$.
\end{remark}

\section{Open analytic problems}
\label{sec:open_problems}

\begin{question}
\label{question:epsilon_distribution}
What is the distribution of
\begin{equation}
 \varepsilon_k
 =\left\lceil\frac{\gamma_k\log p_k}{\pi}\right\rceil
 -\frac{\gamma_k\log p_k}{\pi}
 \label{eq:epsilon_open_distribution}
\end{equation}
in $[0,1)$?  Any equidistribution or discrepancy theorem would immediately
refine the displacement statistics and the prime-power resonance counts.
\end{question}

\begin{conjecture}[Cell equidistribution]
\label{conj:cell_equidistribution}
The sequence $(\varepsilon_k)$ is equidistributed in $[0,1)$.
\end{conjecture}

\begin{remark}
\label{rem:conjecture_consequences}
Conjecture~\ref{conj:cell_equidistribution} would imply, for every continuous
$f$ on $[0,1]$,
\begin{equation}
 \frac1K\sum_{k\leq K}
 f\left(\frac{2(z_k-p_k)}{(\log k)^2}\right)
 \longrightarrow\int_0^1f(t)\dd t,
 \label{eq:displacement_distribution_consequence}
\end{equation}
provided the uniform error in
\eqref{eq:precise_displacement_asymptotic} is incorporated.  It would also
give statistical information on the resonance factors
$\cos(\pi r\varepsilon_k)$.
\end{remark}

\begin{question}
\label{question:marked_trace_formula}
Can the marked measure \eqref{eq:marked_atomic_measure} be realized as the
spectral measure of an operator or as the divisor of a two-variable analytic
function whose arithmetic side is expressed directly through ordinary primes
and prime powers?
\end{question}

\begin{question}
\label{question:regularized_recovery}
Does there exist a regularization of $D_{\ZP}(s,u)$ at or beyond its natural
half-plane that admits stable coefficient extraction in $u$ and Mellin
inversion in $z$, with quantitative control compatible with
\eqref{eq:inverse_stability_bound}?
\end{question}

\begin{question}
\label{question:low_power_resonance}
For fixed small $R$, can one prove nonrandom structure in
\begin{equation}
 \min_{1\leq r\leq R}\|r\varepsilon_k\|_{\R/\mathbb Z}
 \label{eq:low_power_resonance_statistic}
\end{equation}
that correlates with the low prime-power modes in a regularized explicit
formula?
\end{question}

\section{Summary of proved characterizations}
\label{sec:summary}

The complete marked construction is
\begin{equation}
 \boxed{
 n_k=\left\lceil\frac{\gamma_k\log p_k}{\pi}\right\rceil,
 \qquad
 z_k=\exp\left(\frac{\pi n_k}{\gamma_k}\right),
 \qquad
 \mathfrak z_k=(n_k,z_k).
 }
 \label{eq:summary_construction}
\end{equation}
Its exact inverse is
\begin{equation}
 \boxed{
 \gamma_k=\frac{\pi n_k}{\log z_k}.
 }
 \label{eq:summary_inverse}
\end{equation}
The one-cell bounds are
\begin{equation}
 \boxed{
 0\leq\log\frac{z_k}{p_k}<\frac{\pi}{\gamma_k},
 \qquad
 0\leq z_k-p_k<p_k\left(\e^{\pi/\gamma_k}-1\right).
 }
 \label{eq:summary_exact_bounds}
\end{equation}
The principal asymptotics are
\begin{equation}
 \boxed{
 n_k\sim2k,
 \qquad
 z_k\sim p_k\sim k\log k,
 \qquad
 z_k-p_k
 =\left(\frac{\varepsilon_k}{2}+o(1)\right)(\log k)^2.
 }
 \label{eq:summary_asymptotics}
\end{equation}
The counting law is
\begin{equation}
 \boxed{
 \Pi_{\ZP}(x)\sim\pi(x)\sim\frac{x}{\log x}.
 }
 \label{eq:summary_counting}
\end{equation}
The exact matched prime-power resonance is
\begin{equation}
 \boxed{
 \e^{\ii\gamma_k\log p_k^r}
 =(-1)^{rn_k}\e^{-\ii\pi r\varepsilon_k},
 \qquad
 \dist(\gamma_k\log p_k^r,\pi\mathbb Z)
 =\pi\|r\varepsilon_k\|_{\R/\mathbb Z}.
 }
 \label{eq:summary_resonance}
\end{equation}
Finally, the theoretical localization coordinate is
\begin{equation}
 \boxed{
 (n,z)\longmapsto\Gamma(n,z)=\frac{\pi n}{\log z}.
 }
 \label{eq:summary_lattice_map}
\end{equation}
The marked pair is therefore an exact coordinate for an ordinate, while the
analytic problem is to characterize or recover the correct marked support by
an independent global relation.

\begin{thebibliography}{99}

\bibitem{AriasToulisse2013}
J.~Arias de Reyna and J.~Toulisse,
\emph{The $n$-th prime asymptotically},
J. Th\'eor. Nombres Bordeaux \textbf{25} (2013), no.~3, 521--555.
DOI: \href{https://doi.org/10.5802/jtnb.847}{10.5802/jtnb.847}.

\bibitem{Beurling1937}
A.~Beurling,
\emph{Analyse de la loi asymptotique de la distribution des nombres
premiers g\'en\'eralis\'es. I},
Acta Math. \textbf{68} (1937), 255--291.

\bibitem{Diamond1969}
H.~G. Diamond,
\emph{The prime number theorem for Beurling's generalized numbers},
J. Number Theory \textbf{1} (1969), 200--207.
DOI: \href{https://doi.org/10.1016/0022-314X(69)90038-9}{10.1016/0022-314X(69)90038-9}.

\bibitem{Guinand1948}
A.~P. Guinand,
\emph{A summation formula in the theory of prime numbers},
Proc. London Math. Soc. (2) \textbf{50} (1948), 107--119.
DOI: \href{https://doi.org/10.1112/plms/s2-50.2.107}{10.1112/plms/s2-50.2.107}.

\bibitem{Ingham1932}
A.~E. Ingham,
\emph{The Distribution of Prime Numbers},
Cambridge Tracts in Mathematics and Mathematical Physics, no.~30,
Cambridge University Press, Cambridge, 1932.

\bibitem{Titchmarsh1986}
E.~C. Titchmarsh,
\emph{The Theory of the Riemann Zeta-Function},
2nd ed., revised by D.~R. Heath-Brown,
Oxford University Press, Oxford, 1986.

\bibitem{Weil1952}
A.~Weil,
\emph{Sur les ``formules explicites'' de la th\'eorie des nombres premiers},
Comm. S\'em. Math. Univ. Lund, Tome Suppl\'ementaire (1952), 252--265.

\end{thebibliography}

\end{document}
